% Part: first-order-logic
% Chapter: syntax-and-semantics
% Section: introduction

\documentclass[../../../include/open-logic-section]{subfiles}

\begin{document}

\olfileid{fol}{syn}{int}

\olsection{تمهيد}

لا سبيل إلى إقامة نظرية منطق الرتبة الأولى وما وراءها قبل تحديد تركيب
عباراته ودلالاتها. وعبارات هذا المنطق حدود و!!{formula}s. فأما الحدود
فتؤلف من~!!{variable}s و!!{constant}s و!!{function}s. وأما
!!^{formula}s فأولها ما يؤلف من~!!{predicate}s والحدود، وهذه أبسط
!!{formula}s وتسمى «ذرية»؛ ثم تُركّب من~!!{formula}s الذرية صيغ أوسع
باستعمال الروابط المنطقية والأسوار. ولصياغة قواعد التركيب وجوه كثيرة،
وإنما نأخذ منها وجهاً واحداً. فلغيرنا أن يختار رموزاً أخرى، أو يجعل
روابط أخرى أولية، أو يتصرف في الأقواس على غير طريقتنا، أو يستغني عنها
جملة كما في المسمى بالترميز البولندي. غير أن هذه الطرق كلها تجتمع في أن
قواعد التركيب تحد مجموعة الحدود و!!{formula}s حداً \emph{استقرائياً}.
فإذا أُحكم الحد، لم يكن لتوليد كل عبارة بحسب القواعد، في الجوهر، إلا طريق
واحد. ولما كانت العبارة الناشئة عن هذا الحد الاستقرائي~\emph{وحيدة
القراءة}، جاز أن نسند إليها معناها بالطريق نفسه، أي بحد استقرائي.

وباب إسناد المعاني إلى العبارات هو علم الدلالة، وعمدة هذا الباب الإشباع
في~!!a{structure}. فإن~!!^a{structure} تجعل للبنات اللغة معاني:
فـ!!a{domain} مجموعة أشياء غير خالية، والأسوار جارية على عناصر هذا
المجال، ويُجعل لكل~!!{constant}s عنصر منه، ولكل رمز من~!!{function}s
ذي رتبة~$\OLId{latin-ordinary}{n}$ دالة من~$\OLId{latin-looped}{D}^\OLId{latin-ordinary}{n}$ إلى~$\OLId{latin-looped}{D}$، حيث~$\OLId{latin-looped}{D}$ هو~!!{domain}؛ ويُجعل
لـ!!{predicate}s علاقات على~!!{domain}. ومجموع~!!{domain} وهذه المعاني
المسندة إلى المفردات الأصلية هو~!!a{structure}. وقد تقع~!!^{variable}s
في~!!{formula}s، فلا بد لإعطائها دلالة من إسناد~!!{element}s من
~!!{domain} إليها أيضاً، وهو إسناد المتغيرات. ثم تجمع علاقة الإشباع هذه
المعاني: فقد تُشبَع~!!^a{formula} في~!!a{structure}~$\Struct{M}$ بالنسبة
إلى إسناد متغيرات~$\OLId{latin-ordinary}{s}$، ونرمز إلى ذلك بـ$\Sat{\OLId{latin-looped}{M}}{!A}[\OLId{latin-ordinary}{s}]$. وتحد هذه
العلاقة هي الأخرى بالاستقراء على تركيب~$!A$، مستعملين جداول صدق الروابط؛
فنحد، مثلاً، إشباع~$!A \land !B$ بإشباع~$!A$ و$!B$ أو بعدمه. ثم يظهر أن
إسناد المتغيرات لا أثر له إذا كانت~!!{formula}~$!A$ !!a{sentence}، أي
خالصة من المتغيرات الحرة؛ وحينئذ نتكلم بلا قيد على إشباع~!!{sentence}s
أو عدم إشباعها في~!!{structure}s.

ومن علاقة الإشباع~$\Sat{\OLId{latin-looped}{M}}{!A}$ في الجمل نحد المعاني الدلالية الكبرى:
الصحة واللزوم وقابلية الإشباع. فالجملة صحيحة، أي~$\Entails !A$، إذا
أشبعتها كل بنية. وتكون لازمة من مجموعة~!!{sentence}s، أي
$\OLId{greek-variable}{Gamma} \Entails !A$، إذا كان كل~!!{structure} يُشبع جميع~!!{sentence}s
التي في~$\OLId{greek-variable}{Gamma}$ مشبعاً كذلك لـ$!A$. وتكون مجموعة الجمل قابلة للإشباع
إذا وُجدت~!!{structure} تشبع كل~!!{sentence}s فيها معاً. ولأن
~!!{formula}s محدودة استقرائياً، ولأن الإشباع محدود هو أيضاً بالاستقراء
على تركيب~!!{formula}s، جاز لنا أن نستعمل الاستقراء لإثبات أحكام دلالتنا
وبيان الصلات بين المعاني الدلالية المحدودة.




\end{document}
